% 5-4-future-work.tex
%  Budget: ~2pp.  HANDOFF 16.5 -- each recommendation must be grounded in a
%  specific limitation this study actually hit. No "more data, more models".
%  Source map: probabilistic forecasting -> ziel_probabilistic_2018.
%  Rules: decimals in \lr{}; integers Persian; every table cell in \lr{}.

\section{پیشنهادات برای تحقیقات آینده}
\label{sec:5-4}

پیشنهادهای زیر عمومی نیستند. هر کدام از یک محدودیت مشخص برمی‌خیزد که این پژوهش
واقعاً به آن برخورد و در \cref{sec:5-2} نامش را برد.

\subsection*{یک: جداکردنِ دو سازوکارِ درهم‌آمیخته}

مهم‌ترین کارِ نکرده‌ی این پژوهش، اجرای دوباره‌ی آزمون خارج از توزیع با مسیر
سرویس‌دهیِ تصحیح‌شده است. تشخیصِ پشتیبانیِ ویژگی نشان داد که سری بار در آموزش و
در سرویس‌دهی دو تعریف متفاوت داشته است، اما هیچ امتیازدهیِ تازه‌ای انجام نشد،
چون این کار نتایج تثبیت‌شده را دوباره باز می‌کرد. پژوهشی که بار را با تعریف
یکسان تغذیه کند و سپس همان هفت مدل را دوباره امتیاز دهد، برای نخستین بار
\emph{باقیمانده} را می‌سنجد: آن بخش از افت دقت که پس از رفع ناسازگاری برجا
می‌ماند و تنها آن را می‌توان به تغییر بازار نسبت داد. تا پیش از آن اجرا، سهم
بازار کمّی نشده باقی می‌ماند. این پیشنهاد، مشخص‌ترین و کم‌هزینه‌ترین پیشنهاد این
فهرست است: نه مدل تازه‌ای لازم دارد، نه داده‌ی تازه‌ای.

\subsection*{دو: بازآموزی و تصحیح، به‌عنوان یک سیاست مشروط}

\cref{tab:ood-recal} نشان داد تصحیح اریبیِ غلتان به مدل‌هایی که بیشترین انحراف
را داشتند کمک چشمگیری می‌کند و به دو مدلِ کم‌اریب \emph{آسیب} می‌زند. پس یک
سیاستِ «همیشه تصحیح کن» نادرست است. آنچه لازم است، معیاری است که تصمیم بگیرد
تصحیح چه زمانی روشن شود \lr{—} برای مثال بر پایه‌ی پایشِ پیوسته‌ی همان سنجه‌ی
پشتیبانیِ ویژگی، که در این پژوهش تنها یک بار و پس از وقوعِ شکست محاسبه شد.
سنجه‌ای که هر روز محاسبه شود می‌تواند پیش از افت دقت هشدار دهد، چون
بیرون‌رفتنِ ورودی از پشتیبانیِ آموزشی زودتر از بالارفتنِ خطا دیده می‌شود.

\subsection*{سه: گزارش توزیعِ بذر، نه یک اجرای منفرد}

نتیجه‌ی تکمیلیِ \cref{sec:4-5-2} نشان داد که تنها با میانگین‌گیری بر چند بذرِ
تصادفی، \lr{MAE} شبکه‌ی \lr{LSTM} از \lr{3.8734} به \lr{3.6460} می‌رسد. این
تفاوت از فاصله‌ی میان چند مدلِ رقیب در جدول اصلی بزرگ‌تر است. نتیجه‌ی روشن آن
است که یک اجرای تک‌بذری از یک مدل عمیق، پایه‌ی محکمی برای رتبه‌بندی نیست.
پژوهش آینده باید مدل‌های عمیق را با توزیعی از بذرها گزارش کند، نه با یک عدد \lr{—}
و آزمون‌های معناداری را نیز بر همان توزیع بنا کند. کار حاضر برای حفظ انسجام با
نتایج تثبیت‌شده‌ی خود این کار را نکرد و همان انتخاب را صریح اعلام کرد.

\subsection*{چهار: تکرار بر بازارها و پنجره‌های دیگر}

تمام آنچه گزارش شد، بر یک بازار، یک پنجره‌ی آزمون و یک تاریخ تثبیت به دست آمده
است. به‌ویژه یافته‌ی رژیم‌محور به تکرار نیاز دارد: کل بهبود آن بر ۷۷ روز پرتنش
متکی است، و ۷۷ روز برای توانِ آماریِ راحت کم است. تکرار همان روش بر بازاری دیگر
\lr{—} که در این پروژه تنها یک تغییر پیکربندی است \lr{—} یا بر پنجره‌ای طولانی‌تر،
مشخص می‌کند که این یافته خاصیتِ روش است یا خاصیتِ این دو سال. همچنین باید آزموده
شود که آستانه‌ی رژیم، که اینجا یک عدد ثابت از آمارِ دوره‌ی آموزش است، در برابر
تعریفی که خود در طول زمان به‌روز می‌شود چگونه عمل می‌کند.

\subsection*{پنج: مجموعه‌ی ویژگیِ غنی‌تر از آنچه بنچمارک اجازه می‌داد}

\cref{sec:3-4} فهرستی از متغیرهایی را آورد که عمداً در ماتریس ویژگی نیستند، و
برای هر کدام دلیلی متفاوت ذکر کرد. قیمت سوخت به دلیل مجوز تجاری کنار ماند، و
هواشناسی و جریان‌های مرزی به دلیلی که به داده ربطی نداشت: افزودنشان هم‌مقیاسیِ
مقایسه با اعداد منتشرشده را از میان می‌برد. آن مبادله برای این پژوهش درست بود،
چون مقایسه‌پذیری شرطِ لازمِ ارزیابیِ روش بود؛ اما برای پژوهشی که دیگر خود را به
پروتکل مرجع مقید نمی‌داند، درست نیست.

بنابراین مسیر طبیعیِ بعدی، رهاکردنِ آگاهانه‌ی همان قید و افزودن این متغیرهاست:
هواشناسی، جریان‌های مرزی، و \lr{—} در صورت دسترسیِ مجاز به داده \lr{—} شاخص‌های
حاشیه‌ی سوخت. ساختار کد این را از پیش پیش‌بینی کرده است: بلوک‌های ویژگیِ فیزیکی
به‌صورت افزودنی و به‌طور پیش‌فرض خاموش پیاده‌سازی شده‌اند، پس افزودن آن‌ها مسیر
پیش‌فرضی را که نتایج تثبیت‌شده از آن آمده‌اند تغییر نمی‌دهد. نکته‌ی مهم آن است که
چنین پژوهشی دیگر نمی‌تواند نتیجه‌ی خود را مستقیماً در برابر اعداد مرجع بگذارد، و
باید این را صریح بنویسد.

\subsection*{شش: از پیش‌بینی نقطه‌ای به پیش‌بینی احتمالاتی}

هر عددی در این پایان‌نامه یک پیش‌بینی نقطه‌ای است. اما شکستِ خارج از توزیع دقیقاً
از جنسی بود که یک پیش‌بینی نقطه‌ای نمی‌تواند بیانش کند: مدل‌ها با همان اطمینانِ
ظاهریِ همیشه، اعدادی می‌دادند که به‌طور سامان‌مند پایین بودند. پیش‌بینیِ
احتمالاتی \cite{ziel_probabilistic_2018} می‌تواند این وضعیت را به شکل بازه‌هایی
گزارش کند که هنگام دورشدنِ ورودی‌ها از دامنه‌ی آموزش پهن می‌شوند، و برای
تصمیم‌گیرنده‌ای که با خروجی چنین مدلی کار می‌کند، این پهن‌شدن اطلاعاتی است که
یک عدد تنها هرگز نمی‌دهد.
